\begin{abstract}
\noindent
Causal Physics takes the causal order of events, not spacetime geometry, as the primitive of physics.
We formulate it so that every step can be confronted with a pre-registered numerical test, and we
report the outcome of these tests, positive and negative.

\smallskip\noindent
\textbf{Substrate.} The tests decide the nature of the substrate. A discrete substrate with a
finite number of neighbours per event and a rest frame is not supported; the substrate is a
Lorentz-invariant causal set, a Poisson sprinkling in which geometry is read from order and number.
Matter is a flow of events that carries its own four-velocity, and the action of a flow is its proper
time as counted by the order.

\smallskip\noindent
\textbf{Gravity as the response of the order.} Near matter the order changes: clocks slow down, causal
cones narrow and events become denser. A fixed order gives the scalar (Nordstr\"om) class
$\gamma=-1$, cones that only narrow give $\gamma=1/d_s$, and an order that conserves its causal
throughput gives $\gamma=2/(d_s-1)$, which is $\gamma=1$ exactly in three spatial dimensions. In
3+1 dimensions this throughput principle is the linearised coupling of general relativity to the full
energy--momentum tensor. A substrate rule in which every flow narrows the cone in its own rest
frame and capacity deficits add reproduces this coupling on 3+1 sprinklings: two counter-streaming
beams gravitate with the predicted active mass ($2.40\pm0.05$ against $2.36$, alternatives $1.56$
and below), the cones become anisotropic and a single beam drags them.

\smallskip\noindent
\textbf{The rule from the action of the order.} With proper time counted by the order, the substrate
rule is the linear response of an order whose action has the Einstein--Hilbert stiffness. We compute
the expectation value of the smeared Benincasa--Dowker action of a causal set to second order in a
static perturbation, in 3+1, by deterministic quadrature with nothing from general relativity put in.
The Einstein--Hilbert stiffness emerges in the conformal sector, $6.0003$ against $6$, and in the mixed
sector, where the perturbation changes the causal order itself, $7.983$ against $8$, and in the optical
sector, where rays bend, $2.002$ against $2$. With all three coefficients $\gamma=1$ follows from the action of
the order ($0.997$ to $1.000$). The discreteness corrects the stiffness non-analytically, $1-\hat K\simeq(k\ell_k)^2[4.0\ln(1/k\ell_k)+0.9]$, and adds a
non-local term that grows with the age $T$ of the static field, $-0.10\,(k\ell_k)^3T/\ell_k$, which
we predicted analytically before computing it.

\smallskip\noindent
\textbf{Predictions.} Beyond the leading order the sectors differ, and each weakens with the age $T$ of a
static field: linearly in the conformal and mixed sectors, faster than linearly in the optical sector within
the computed range. The age terms are not suppressed like the local corrections: if the age of a static field
is cosmological, a non-locality scale at the level of existing bounds, $\ell_k\sim10^{-19}$\,m, places the
breakdown of the static response at tens of micrometres or above, within reach of torsion balances. The
asymptotic form of the optical age term, which fixes the far-field law of a point mass, is the next
computation. Beyond the linear order the theory is completed by general relativity, to be
derived thermodynamically; strong-field predictions of a scalar field equation are therefore not
predictions of the theory. Every result quoted here is reproducible with the open Causality
Lab~\cite{CausalityLab}, including the negative ones.
\end{abstract}
